\section{What to Write Down, and Where}
\label{sec:clocks-write-down}

A clock fact has three places it can live, and only one of them survives the project.

\subsection{The Three Places}

\begin{table}[htbp]
    \centering
    \footnotesize
    \begin{tabular}{@{}p{0.28\textwidth}p{0.30\textwidth}p{0.34\textwidth}@{}}
        \toprule
        \textbf{Place} & \textbf{Lifetime} & \textbf{What happens to it} \\
        \midrule
        In the chat where you found it & minutes & lost \\
        In a commit message & months, if you search & nobody searches \\
        In the source, next to the constant it explains & the life of the code & \textbf{read by the next person} \\
        \bottomrule
    \end{tabular}
    \caption{The only durable location is the one somebody will read while editing the thing it
    describes.}
    \label{tab:where-facts-live}
\end{table}

\subsection{How the Numbers Here Are Written}

Every measured constant in this project carries three things: the number, the unit, and how it was
obtained. For example:

\begin{minted}{rust}
/// Half a bit period on the bit-banged I2C bus, in loop iterations.
///
/// Measured: 60 iterations take ~5 us at SYSCLK = 19 MHz on this board, so ~95 kHz.
/// This is a measurement, not a datasheet value -- `asm::delay` is about three cycles
/// per iteration here, so the divisor is not the clock.
const HALF_BIT_CYCLES: u32 = 60;
\end{minted}

and for the value that everything else depends on:

\begin{minted}{rust}
/// SYSCLK as measured with a stopwatch over 20 blink cycles: 19 MHz (the HFRCO default).
/// The part supports 78 MHz; this board's firmware never selects it.
pub const DEFAULT_SYSCLK_HZ: u32 = 19_000_000;
\end{minted}

\subsection{The Table That Belongs in Every Project}

\begin{table}[htbp]
    \centering
    \footnotesize
    \begin{tabular}{@{}p{0.28\textwidth}p{0.24\textwidth}p{0.40\textwidth}@{}}
        \toprule
        \textbf{Quantity} & \textbf{Value here} & \textbf{How it was measured} \\
        \midrule
        SYSCLK & 19 MHz & stopwatch over 20 blink cycles \\
        \codeinline{asm::delay} cost & ~3 cycles/iteration & cycle counter, two build profiles \\
        I2C half period & ~5 $\mu$s (60 iterations) & scope-free: 95 kHz against a known device \\
        Buzzer tone & 2.7 kHz, TOP = 3518 & TIMER0 TOP read back over SWD \\
        IMU sample rate & 416 Hz configured & \codeinline{CTRL1\_XL} register + measured intervals \\
        Sampling loop & 158 Hz when claiming 400 & cycle counter, per-sample interval recorded \\
        \bottomrule
    \end{tabular}
    \caption{The clock facts of this project. The last line is the one that mattered: a loop that
    reports a rate it does not achieve poisons every timing test downstream.}
    \label{tab:clock-facts}
\end{table}

\subsection{The One That Generalises}

The last row of that table is the lesson of this chapter. The firmware said 400 Hz because the
\emph{intended} interval was 2500 $\mu$s. It ran at 158 Hz because the display flush blocked for tens
of milliseconds inside the loop. Every threshold in the state machine was therefore 2.5× fast, and
three rounds of drop tests were interpreted against a wrong time base before anyone noticed that four
different drops showed the same 0.62 s of calm.

The fix was not a better delay. It was to \textbf{measure the interval between samples and pass the
measurement on}:

\begin{minted}{rust}
let now = dwt::cycles();
let dt_us = dwt::cycles_to_us(dwt::since(self.last), sysclk_hz);
self.last = now;
// The classifier is told what actually happened, not what was planned.
detector.update(&accel, dt_us);
\end{minted}

Every timing defect in this project — and there were four — would have been caught on the day it was
introduced by that three-line habit.
